% Part: second-order-logic
% Chapter: syntax-and-semantics
% Section: language-of-sol

\documentclass[../../../include/open-logic-section]{subfiles}

\begin{document}

\olfileid{sol}{syn}{lan}

\olsection{لغة منطق الرتبة الثانية}

كما في منطق الرتبة الأولى، تُبنى تعبيرات منطق الرتبة الثانية
من مفردات أساسية تضم \emph{!!{variable}s} و\emph{!!{constant}s}
و\emph{!!{predicate}s}، وتضم أحيانًا \emph{!!{function}s}. ومن هذه المفردات،
مع الروابط المنطقية والمكمّمات وعلامات الترقيم كالأقواس والفواصل،
تُنشأ \emph{الحدود} و\emph{!!{formula}s}. والفرق هو أن منطق الرتبة الثانية،
إضافة إلى متغيرات الأفراد، يشتمل أيضًا على متغيرات للعلاقات والدوال،
ويتيح التكميم عليها. ولذلك تكون رموزه المنطقية هي رموز منطق الرتبة
الأولى، مضافًا إليها ما يأتي:

\begin{enumerate}
\item مجموعة !!{denumerable} من !!{variable}s العلاقات من الرتبة الثانية،
  من كل رتبة~$n$: $\Obj V_0^n$، $\Obj V_1^n$، $\Obj V_2^n$، \dots
\item مجموعة !!{denumerable} من !!{variable}s الدوال من الرتبة الثانية،
  من كل رتبة~$n$: $\Obj u_0^n$، $\Obj u_1^n$، $\Obj u_2^n$، \dots
\end{enumerate}

في منطق الرتبة الأولى، تُدرج عادةً علاقة !!{identity}~$\eq$. وفي منطق
الرتبة الأولى أيضًا، تعد الرموز غير المنطقية في لغة~$\Lang{L}$
ضرورية لكي نستطيع التعبير عن أمور ذات شأن. وثمة بالطبع !!{sentence}s
لا تستخدم أي رمز غير منطقي، غير أنه يصعب، باستعمال~$\eq$ وحدها، قول
شيء ذي شأن. أما في منطق الرتبة الثانية، فلأن لدينا عددًا غير محدود
من متغيرات العلاقات والدوال، يمكننا قول كل ما يمكن قوله في لغة من
الرتبة الأولى حتى من دون مخزون خاص من الرموز غير المنطقية.

\end{document}
